\begin{onehalfspace}

Questa sezione approfondisce le Sezioni~\ref{sec:surface_concetto}--\ref{sec:oltre_pauli}
del Capitolo~\ref{chap:risultati}. Nel corpo restano curve, soglie, legge
di scala e conclusioni; qui sono conservati il funzionamento del codice,
la griglia sperimentale, le verifiche di riproducibilità e il controllo
sui canali non-Pauli.

\subsection{Surface code e decodifica MWPM}
\label{app:surface_teoria}

Il surface code~\cite{kitaev2003,fowler2012} dispone su una griglia
bidimensionale qubit dati e qubit di misura. Gli stabilizzatori sono locali
e un patch planare ruotato di distanza $d$ usa $d^2$ dati, oltre a un
numero dello stesso ordine di ancille. Le misure sono ripetute per più
round perché anche la sindrome è rumorosa; una variazione fra due cicli
costituisce un \emph{detection event}.

Nel bulk, una catena d'errore produce tipicamente detection events in
coppie; al bordo uno dei due estremi può terminare sul confine. Il
\ac{MWPM} associa gli eventi minimizzando un peso collegato alla
probabilità delle catene~\cite{pymatching2021}. Se errore reale e
correzione differiscono per un ciclo che attraversa il patch, si verifica
un errore logico silenzioso: la sua frequenza è $p_L$.

La stima richiede milioni di circuiti con centinaia di qubit, ma i memory
experiment sono circuiti stabilizer. Stim~\cite{stim2021} sfrutta il
teorema di Gottesman--Knill per campionare detection events; PyMatching
costruisce il grafo dal \ac{DEM} e decodifica. Il Listato~\ref{lst:stim_pymatching}
nella Sezione~\ref{app:codice} documenta l'implementazione. Shor non è un
circuito stabilizer: per questo i tassi del memory experiment non sono
trasferiti direttamente alla successiva analisi di sensibilità.

\subsection{Griglia sperimentale e riproducibilità}
\label{app:surface_protocollo}

\begin{table}[H]
\centering
\begin{tabularx}{\textwidth}{llX}
\toprule
\textbf{Variabile} & \textbf{Valori} & \textbf{Osservazione attesa} \\
\midrule
distanza $d$ & 3, 5, 7, 9 & $p_L$ diminuisce con $d$ sotto soglia \\
rounds & $d$ e $2d$ & cicli di memoria quantistica \\
$p$ fisico & $10^{-4}$ \dots $2\times10^{-2}$ & curva $p$--$p_L$ \\
shot & $\geq10^4$ per punto & intervalli di confidenza stabili \\
base & memory X e memory Z & controllo delle due protezioni \\
\bottomrule
\end{tabularx}
\caption[Disegno sperimentale surface code]{Griglia prescritta. La
campagna eseguita copre le quattro distanze, entrambe le basi,
$\mathrm{rounds}=d$ e $p\in[2\times10^{-3},10^{-2}]$; lo scaling a
$2d$ resta un'estensione.}
\label{tab:surface_esperimento}
\end{table}

Il modello è circuit-level: depolarizzazione dopo ogni Clifford e flip di
reset e misura con probabilità $p$. Il campione è dimensionato punto per
punto per raccogliere circa duecento eventi logici, da $2\times10^5$ a
$4\times10^7$ shot. Una prima passata da $2\times10^5$ stima l'ordine di
grandezza di $p_L$; la seconda fissa la numerosità. A $d=9$ e
$p=2\times10^{-3}$ una numerosità fissa darebbe due soli eventi e il
$70\%$ di incertezza relativa, mentre il campionamento adattivo la porta
al $3.5\%$. Seed e numerosità sono registrati per ogni punto. Una
ripetizione indipendente ha riprodotto tutte le stime entro l'errore
statistico e le stesse soglie.

\subsection{Errori coerenti, leakage e controllo full-state}
\label{app:surface_non_pauli}

Stim rappresenta gate Clifford ed errori Pauli casuali. L'hardware può
però produrre sovrarotazioni coerenti, leakage verso $\ket{2}$ e
correlazioni temporali; una simulazione solo Clifford--Pauli può quindi
essere ottimistica~\cite{plaquette2026}. In una sovrarotazione
$U=e^{-i\varepsilon H}$ le ampiezze si sommano coerentemente, mentre un
errore stocastico accumula probabilità.

Il divario è illustrato da un qubit soggetto a $N$ operazioni con
$p=10^{-2}$, modellate come X stocastico oppure come
$R_X(\theta)$ con $\sin^2(\theta/2)=p$. Le probabilità esatte sono
$[1-(1-2p)^N]/2$ e $\sin^2(N\theta/2)$; nel tratto iniziale crescono
come $Np$ e $N^2p$. A $N=8$ il secondo errore è circa sette volte il
primo per i parametri scelti.

\begin{figure}[H]
\centering
\includegraphics[width=0.75\textwidth]{qec_coherent.pdf}
\caption[Errore coerente vs errore di Pauli]{Accumulo su un qubit a
parità di probabilità marginale per operazione. La simulazione full-state
(marcatori) coincide con le espressioni teoriche; il fattore osservato non
è universale.}
\label{fig:qec_coherent}
\end{figure}

Lo studio di riferimento riporta, in altri regimi, errori logici
$1.7\%$ contro $23.5\%$ in presenza di leakage e $18.1\%$ contro
$38.7\%$ con errori coerenti; per un surface code a distanza 9 il divario
può raggiungere un fattore 25--56 e spostare la soglia di circa il
$20\%$~\cite{plaquette2026}. Questi valori non sono trasferiti alla
campagna della tesi: motivano un approccio channel-first, nel quale il
canale fisico è definito una volta e approssimato solo dopo aver stabilito
quale informazione preservare.

Il controllo interno usa un codice a distanza 3 simulato full-state per
quattro cicli, confrontando X di Pauli e $R_X$ coerente a pari
probabilità marginale. A $p=3\times10^{-2}$ il canale coerente produce
meno flip grezzi ($0.237$ contro $0.262$) ma un errore MWPM leggermente
maggiore ($0.0789$ contro $0.0767$, $+2.9\%$). Le misure ripetute
interrompono parte dell'accumulo coerente, senza equivalere a un twirling
Pauli completo.

Il decoder ibrido ottiene guadagni $1.019$, $1.036$, $1.041$ sotto Pauli
e $1.029$, $1.048$, $1.052$ sotto sovrarotazione, per
$p=1,2,3\times10^{-2}$. Lo scarto è circa una deviazione standard e non
è significativo. Il risultato esclude quindi la regola semplicistica
``non-Pauli implica vantaggio appreso'': ciò che il matching non
rappresenta è la non decomponibilità in archi, come nel crosstalk che
accende più di due detector. Le soglie del corpo restano baseline entro
il modello Pauli; leakage e canali più ricchi sono sviluppi futuri.

\end{onehalfspace}
